\pagebreak

\section{Common Xilinx Procedures Library File - xilinx/common.3pl}

    \index{library!xilinx/common.3pl}
    This library provides common modules, procedures and functions for
    all Xilinx FPGAs.
    
    This library must be incorporated by using a {\bf source} statement -
    \begin{lstlisting}[gobble=4, language=threepl]
       source "xilinx/common.3pl"
    \end{lstlisting}
    
    {\bf All Xilinx family header files include this library.}



    \subsection{library modules}
    


        \subsubsection{sequence - generate a sequence}\label{sec:lmsequence}
            
            {\bf Library: xilinx/common.3pl}

            {\bf Prototype:} \\
            \vsp
            \begin{lstlisting}[gobble=12, language=threepl]
               global module
               sequence (value(out) <- var(s), value(step, "log", true)) {}
            \end{lstlisting}

            {\bf Output parameters:} \\
            {\bf out} is the output value variable.

            {\bf Input parameters:} \\
            {\bf s} is the sequence array.\\
            {\bf step} is an optional input to control data movement.

            {\bf description:} \\
            This module generates a cyclic sequence of values, these
            appearing on the output value variable. The sequence is
            determined by the input argument {\bf s} which must be an
            array of immediate values. The type of these values must be
            an immediate equivalent to the type of the output argument
            and may be of compound type. The 2nd input argument, {\bf step},
            is optional and is used to control when the sequence advances,
            otherwise this occurs on every clock edge.




    \subsection{library procedures}

        
        
        \subsubsection{IOconstraint - add an I/O constraint to the constraints file}\label{sec:lpioconstraint}

            
            {\bf Library: xilinx/common.3pl}

            {\bf Prototype:} \\
            \vsp
            \begin{lstlisting}[gobble=12, language=threepl]
               global procedure
               IOconstraint (str(loc), str(con), var(val)) {}
            \end{lstlisting}

            {\bf Output parameters:} \\
            none

            {\bf Input parameters:} \\            
            {\bf loc} is the I/O pad designation \\
            {\bf con} is the constraint identifier \\
            {\bf val} is the constraint value \\
        
            {\bf In-line target execution:} \\
            no

            {\bf description:} \\
            This procedure adds a line to the NCF (ISE) or XDC (Vivado) constraint
            file as determined by directive "postProcessTool" (\autorefph{sec:directives}).
            
            {\bf con} is a string sepcifying the constraint. Valid constraints are -

            \begin{tabular}{| l | l | l |}
            \hline
            "diff\_term":           & log & input differential termination \\ \hline
            "drive":                & int & output drive current           \\ \hline
            "ibuf\_delay\_value":   & int & input direct delay             \\ \hline
            "ifd\_delay\_value":    & int & input flip-flop delay          \\ \hline
            "iostandard":           & log & I/O standard                   \\ \hline
            "pullup":               & log & pullup                         \\ \hline
            "pulldown":             & log & pulldawn                       \\ \hline
            "slew":                 & str & output slew rate               \\ \hline
            "tnm":                  & str & timing name                    \\ \hline
            \end{tabular}
            
            The constraint value {\bf val} may be type "str", "log" or "int"
            as appropriate.

            Note that if the pad is bi-directional and an I/O standard constraint
            is provided separately for both input and output the second will be
            ignored.



        \subsubsection{addncfline - add line to NCF file}\label{sec:lpaddncfline}
            \index{constraint}
            
            {\bf Library: xilinx/common.3pl}

            {\bf Prototype:} \\
            \vsp
            \begin{lstlisting}[gobble=12, language=threepl]
               global procedure
               addncfline (str(s)) {}
            \end{lstlisting}

            {\bf Output parameters:} \\
            none

            {\bf Input parameters:} \\            
            {\bf s} is the line to be added to the NCF file
        
            {\bf In-line target execution:} \\
            no

            {\bf description:} \\
            This procedure adds a line to the NCF file.
            (Note that a trailing semicolon is not appended.)

            See also - \\
            {\bf append\_ncf\_file()} \autorefph{sec:lpappncf}\\
            {\bf elementaddncf()} \autorefph{sec:lpelementaddncf}\\
            {\bf addxdcline()} \autorefph{sec:lpaddxdcline}\\
            {\bf maxdelay()} \autorefph{sec:lpmaxdelay}\\
            {\bf period\_mul()} \autorefph{sec:lpperiodmul}\\
            {\bf period\_div()} \autorefph{sec:lpperiodmul}\\
            {\bf netaddncf()} \autorefph{sec:lpnetaddncf}.



        \subsubsection{addxdcline - add line to XDC file}\label{sec:lpaddxdcline}
            \index{constraint}
            
            {\bf Library: xilinx/common.3pl}

            {\bf Prototype:} \\
            \vsp
            \begin{lstlisting}[gobble=12, language=threepl]
               global procedure
               addxdcline (str(fmt), value(sig)) {}
            \end{lstlisting}

            {\bf Output parameters:} \\
            none

            {\bf Input parameters:} \\            
            {\bf fmt} is a formta string of the line to be added to the XDC file.
            {\bf sig} is an optional signal to insert into the line.
        
            {\bf In-line target execution:} \\
            no

            {\bf description:} \\
            This procedure adds a line to the XDC file.
            (Note that a trailing semicolon is not appended.)
            {\bf fmt} is a format string which may contain a \verb'%n'
            conversion to insert the optional signal identifier {\bf sig}.

            See also - \\
            {\bf addncfline()} \autorefph{sec:lpaddncfline}\\
            {\bf append\_ncf\_file()} \autorefph{sec:lpappncf}\\
            {\bf elementaddncf()} \autorefph{sec:lpelementaddncf}\\
            {\bf maxdelay()} \autorefph{sec:lpmaxdelay}\\
            {\bf period\_mul()} \autorefph{sec:lpperiodmul}\\
            {\bf period\_div()} \autorefph{sec:lpperiodmul}\\
            {\bf netaddncf()} \autorefph{sec:lpnetaddncf}.










        \subsubsection{append\_ncf\_file - append a file to the NCF file}\label{sec:lpappncf}
            
            {\bf Library: xilinx/common.3pl}

            {\bf Prototype:} \\
            \vsp
            \begin{lstlisting}[gobble=12, language=threepl]
               global procedure
               append_ncf_file (str(filename)) {}
            \end{lstlisting}

            {\bf Output parameters:} \\
            none

            {\bf Input parameters:} \\            
            {\bf filename} is the input file name.
        
            {\bf In-line target execution:} \\
            no

            {\bf description:} \\
            This procedure reads a file and appends it to the NCF file
            (Netlist Constraints File).

            See also - \\   
            {\bf addncfline()} \autorefph{sec:lpaddncfline}\\
            {\bf elementaddncf()} \autorefph{sec:lpelementaddncf}\\
            {\bf addxdcline()} \autorefph{sec:lpaddxdcline}\\
            {\bf maxdelay()} \autorefph{sec:lpmaxdelay}\\
            {\bf period\_mul()} \autorefph{sec:lpperiodmul}\\
            {\bf period\_div()} \autorefph{sec:lpperiodmul}\\
            {\bf netaddncf()} \autorefph{sec:lpnetaddncf}.






    


        \subsubsection{elementaddncf - add line to NCF file}\label{sec:lpelementaddncf}
            
            {\bf Library: xilinx/common.3pl}

            {\bf Prototype:} \\
            \vsp
            \begin{lstlisting}[gobble=12, language=threepl]
               global procedure
               elementaddncf (str({id, attr, val})) {}
            \end{lstlisting}

            {\bf Output parameters:} \\
            none

            {\bf Input parameters:} \\
            {\bf id} is the netlist block identifier (optional). \\
            {\bf attr} is the attribute or constraint name or else a line to be added to the NCF file for the element. \\
            {\bf val} is the value of the attribute or constraint if the previous argument
            is an attribute or constraint name. \\
        
            {\bf In-line target execution:} \\
            no

            {\bf description:} \\
            This procedure adds a line to the Xilinx NCF file for an element. If
            the block identifier is given then the NCF entry applies to that
            element, otherwise it applies to the element most recently
            constructed by inbuilt procedure {\bf element()}. Note that in either case for the
            element concerned its construction must have been completed by
            calling inbuilt procedure {\bf elementend()} before calling the
            current library procedure.

            The constraint string will have something like
            \verb'"INST block234"' prepended, where "block234" is the identifier
            for the element, and will also have a trailing ";" appended.
            
            If the argument to {\bf val} is omitted then the string {\bf attr} will
            be used verbatim. If the argument to {\bf val} is present then
            {\bf attr} will be used as a key to a table of constraints for the element
            and {\bf val} will be the associated value. In that case when the attribute
            or constraintnis written to the constraint file a "=" will be inserted between
            the attribute/constraint name and the key. Thus an attribute may be
            entered by -
            \begin{lstlisting}[gobble=12, language=threepl]
               elementaddncf("b123", "REG=3");
            \end{lstlisting}
            in which case the entire second argument becomes the key and the
            value is null. The key only is copied to the constraint file. The
            limitation of this is that the attribute value cannot be changed.
            If on the other hand the form
            \begin{lstlisting}[gobble=12, language=threepl]
               elementaddncf("b123", "REG", "3");
            \end{lstlisting}
            is used then the attribute can be changed later if desired.

            See also - \\       
            {\bf addncfline()} \autorefph{sec:lpaddncfline}\\
            {\bf append\_ncf\_file()} \autorefph{sec:lpappncf}\\
            {\bf addxdcline()} \autorefph{sec:lpaddxdcline}\\
            {\bf maxdelay()} \autorefph{sec:lpmaxdelay}\\
            {\bf period\_mul()} \autorefph{sec:lpperiodmul}\\
            {\bf period\_div()} \autorefph{sec:lpperiodmul}\\
            {\bf netaddncf()} \autorefph{sec:lpnetaddncf}.






        \subsubsection{elementBinProperty - add a property of type "int" to an element in binary format}\label{sec:lpelementBinProperty}
            
            {\bf Library: xilinx/common.3pl}

            {\bf Prototype:} \\
            \vsp
            \begin{lstlisting}[gobble=12, language=threepl]
               global procedure
               elementBinProperty (str(id), str(prop), int({n, min, max}), int(digits)) {}
            \end{lstlisting}

            {\bf Output parameters:} \\
            none

            {\bf Input parameters:} \\            
            {\bf id} is the element identifier or null.\\
            {\bf prop} is the property identifier.\\
            {\bf n} is the property value.\\
            {\bf min} is the property minimum value.\\
            {\bf max} is the property maximum value.\\
            {\bf digits} is the number of binary digits required.\\
        
            {\bf In-line target execution:} \\
            no

            {\bf description:} \\
            This procedure adds a property of type "int" to an element in
            binary format.
            {\bf id} is a unique identifier for the element. If the argument is
            omitted then the current element is used, i.e. the one most recently
            created using inbuilt procedure {\bf element()} \autorefph{sec:ipelement}.
            
            {\bf prop} is the property identifier.
            
            {\bf n} is the value of the property. If no argument is passed then
            this procedure does nothing and returns. Note that if this procedure
            is called within another module or procedure and the property
            value is a parameter of that enclosing procedure then the inbuilt
            fuction {\bf nullarg()} \autorefph{sec:ifnullarg} may be used to enclose the
            argument to this procedure so that an omitted argument is ignored rather
            than the default parameter value of the enclosing procedure being used.
            
            {\bf min} and {\bf max} are the minimum and maximum allowed values
            of the property. If one of these is exceeded then a fatal error
            will be forced.
            
            {\bf digits} specifies the number of binary digits with which to
            express the property value. The binary digit string is not
            preceded by any base indicator (e.g. there is {\bf no} leading "0b").
            
            See also - \\
            {\bf elementLogProperty()} \autorefph{sec:lpelementLogProperty} for logical properties,
            {\bf elementIntProperty()} \autorefph{sec:lpelementIntProperty} for integer properties,
            %{\bf elementBinProperty()} \autorefph{sec:lpelementBinProperty} for binary properties,
            {\bf elementHexProperty()} \autorefph{sec:lpelementHexProperty} for hexadecimal properties,
            {\bf elementFloatProperty()} \autorefph{sec:lpelementFloatProperty} for floating point properties,
            {\bf elementStrProperty()} \autorefph{sec:lpelementStrProperty} for string properties.
        
        


        \subsubsection{elementcheck - check that a Xilinx element is allowed in the current device family}\label{sec:lpelementcheck}
            
            {\bf Library: xilinx/common.3pl}

            {\bf Prototype:} \\
            \vsp
            \begin{lstlisting}[gobble=12, language=threepl]
               global procedure
               elementCheck(str({element, name})) {}
            \end{lstlisting}

            {\bf Output parameters:} \\
            none

            {\bf Input parameters:} \\            
            {\bf element} is the name of a Xilinx library element. \\
            {\bf name} is the name of a library module, procedure or function which is
            attempting to create the library element.
        
            {\bf In-line target execution:} \\
            no

            {\bf description:} \\
            This procedure checks that the Xilinx library element named by the first
            argument is available for the current Xilinx device family. Allowed elements
            for each device family are kept in an abbreviated map in xilinx/common.3pl,
            i.e. only a subset of elements are included, namely those used explicitly
            in 3PL libraries. The map can be expanded when required.
            
            If the check fails, i.e. the element is not available in the current device family,
            then 3PL exits with an error message and traceback. {\bf name} is provided solely
            for use in the error message.
            
            See also library function {\bf elementcheck()} \autorefph{sec:lfelementcheck}
            which can check non-fatally.




        \subsubsection{elementFloatProperty - add a property of type "float" to an element}\label{sec:lpelementFloatProperty}
            
            {\bf Library: xilinx/common.3pl}

            {\bf Prototype:} \\
            \vsp
            \begin{lstlisting}[gobble=12, language=threepl]
               global procedure
               elementFloatProperty (str(id), str(prop), float({n, min, max}), int(digits)) {}
            \end{lstlisting}

            {\bf Output parameters:} \\
            none

            {\bf Input parameters:} \\            
            {\bf id} is the element identifier or null.\\
            {\bf prop} is the property identifier.\\
            {\bf n} is the property value.\\
            {\bf min} is the property minimum value.\\
            {\bf max} is the property maximum value.\\
            {\bf digits} is the number of digits required after the decimal point.\\
        
            {\bf In-line target execution:} \\
            no

            {\bf description:} \\
            This procedure adds a property of type "float" to an element.
            {\bf id} is a unique identifier for the element. If the argument is
            omitted then the current element is used, i.e. the one most recently
            created using inbuilt procedure {\bf element()} \autorefph{sec:ipelement}.
            
            {\bf prop} is the property identifier.
            
            {\bf n} is the value of the property. If no argument is passed then
            this procedure does nothing and returns. Note that if this procedure
            is called within another module or procedure and the property
            value is a parameter of that enclosing procedure then the inbuilt
            fuction {\bf nullarg()} \autorefph{sec:ifnullarg} may be used to enclose the
            argument to this procedure so that an omitted argument is ignored rather
            than the default parameter value of the enclosing procedure being used.
            
            {\bf min} and {\bf max} are the minimum and maximum allowed values
            of the property. If one of these is exceeded then a fatal error
            will be forced.
            
            {\bf digits} specifies the number of digits to appear to the right
            of the decimal point.
            
            See also - \\
            {\bf elementLogProperty()} \autorefph{sec:lpelementLogProperty} for logical properties,
            {\bf elementIntProperty()} \autorefph{sec:lpelementIntProperty} for integer properties,
            {\bf elementBinProperty()} \autorefph{sec:lpelementBinProperty} for binary properties,
            {\bf elementHexProperty()} \autorefph{sec:lpelementHexProperty} for hexadecimal properties,
            %{\bf elementFloatProperty()} \autorefph{sec:lpelementFloatProperty} for floating point properties,
            {\bf elementStrProperty()} \autorefph{sec:lpelementStrProperty} for string properties.







        \subsubsection{elementHexProperty - add a property of type "int" to an element in hexadecimal format}\label{sec:lpelementHexProperty}
            
            {\bf Library: xilinx/common.3pl}

            {\bf Prototype:} \\
            \vsp
            \begin{lstlisting}[gobble=12, language=threepl]
               global procedure
               elementHexProperty (str(id), str(prop), int({n, min, max}), int(digits)) {}
            \end{lstlisting}

            {\bf Output parameters:} \\
            none

            {\bf Input parameters:} \\            
            {\bf id} is the element identifier or null.\\
            {\bf prop} is the property identifier.\\
            {\bf n} is the property value.\\
            {\bf min} is the property minimum value.\\
            {\bf max} is the property maximum value.\\
            {\bf digits} is the number of hexadecimal digits required.\\
        
            {\bf In-line target execution:} \\
            no

            {\bf description:} \\
            This procedure adds a property of type "int" to an element in
            binary format.
            {\bf id} is a unique identifier for the element. If the argument is
            omitted then the current element is used, i.e. the one most recently
            created using inbuilt procedure {\bf element()} \autorefph{sec:ipelement}.
            
            {\bf prop} is the property identifier.
            
            {\bf n} is the value of the property. If no argument is passed then
            this procedure does nothing and returns. Note that if this procedure
            is called within another module or procedure and the property
            value is a parameter of that enclosing procedure then the inbuilt
            fuction {\bf nullarg()} \autorefph{sec:ifnullarg} may be used to enclose the
            argument to this procedure so that an omitted argument is ignored rather
            than the default parameter value of the enclosing procedure being used.
            
            {\bf min} and {\bf max} are the minimum and maximum allowed values
            of the property. If one of these is exceeded then a fatal error
            will be forced.
            
            {\bf digits} specifies the number of hexadecimal digits with which to
            express the property value. The hexadecimal digit string is not
            preceded by any base indicator (e.g. there is {\bf no} leading "0x").
            
            See also - \\
            {\bf elementLogProperty()} \autorefph{sec:lpelementLogProperty} for logical properties,
            {\bf elementIntProperty()} \autorefph{sec:lpelementIntProperty} for integer properties,
            {\bf elementBinProperty()} \autorefph{sec:lpelementBinProperty} for binary properties,
            %{\bf elementHexProperty()} \autorefph{sec:lpelementHexProperty} for hexadecimal properties,
            {\bf elementFloatProperty()} \autorefph{sec:lpelementFloatProperty} for floating point properties,
            {\bf elementStrProperty()} \autorefph{sec:lpelementStrProperty} for string properties.








        \subsubsection{elementIntProperty - add a property of type "int" to an element}\label{sec:lpelementIntProperty}
            
            {\bf Library: xilinx/common.3pl}

            {\bf Prototype:} \\
            \vsp
            \begin{lstlisting}[gobble=12, language=threepl]
               global procedure
               elementIntProperty (str(id), str(prop), int({n, min, max})) {}
            \end{lstlisting}

            {\bf Output parameters:} \\
            none

            {\bf Input parameters:} \\            
            {\bf id} is the element identifier or null.\\
            {\bf prop} is the property identifier.\\
            {\bf n} is the property value.\\
            {\bf min} is the property minimum value.\\
            {\bf max} is the property maximum value.\\
        
            {\bf In-line target execution:} \\
            no

            {\bf description:} \\
            This procedure adds a property of type "int" to an element. {\bf
            id} is a unique identifier for the element. If the argument is
            omitted then the current element is used, i.e. the one most recently
            created using inbuilt procedure {\bf element()}
            \autorefph{sec:ipelement}.
            
            {\bf prop} is the property identifier.

            {\bf n} is the value of the property. If no argument is passed then
            this procedure does nothing and returns. Note that if this procedure
            is called within another module or procedure and the property value
            is a parameter of that enclosing procedure then the inbuilt fuction
            {\bf nullarg()} \autorefph{sec:ifnullarg} may be used to enclose the
            argument to this procedure so that an omitted argument is ignored
            rather than the default parameter value of the enclosing procedure
            being used.

            {\bf min} and {\bf max} are the minimum and maximum allowed values
            of the property. If one of these is exceeded then a fatal error
            will be forced.
            
            See also - \\
            {\bf elementLogProperty()} \autorefph{sec:lpelementLogProperty} for logical properties,
            %{\bf elementIntProperty()} \autorefph{sec:lpelementIntProperty} for integer properties,
            {\bf elementBinProperty()} \autorefph{sec:lpelementBinProperty} for binary properties,
            {\bf elementHexProperty()} \autorefph{sec:lpelementHexProperty} for hexadecimal properties,
            {\bf elementFloatProperty()} \autorefph{sec:lpelementFloatProperty} for floating point properties,
            {\bf elementStrProperty()} \autorefph{sec:lpelementStrProperty} for string properties.
         





        \subsubsection{elementLogProperty - add a property of type "log" to an element}\label{sec:lpelementLogProperty}
            
            {\bf Library: xilinx/common.3pl}

            {\bf Prototype:} \\
            \vsp
            \begin{lstlisting}[gobble=12, language=threepl]
               global procedure
               elementLogProperty (str(id), str(prop), log(l)) {}
            \end{lstlisting}

            {\bf Output parameters:} \\
            none

            {\bf Input parameters:} \\            
            {\bf id} is the element identifier or null.\\
            {\bf prop} is the property identifier.\\
            {\bf l} is the property value.\\
        
            {\bf In-line target execution:} \\
            no

            {\bf description:} \\
            This procedure adds a property of type "log" to an element.
            {\bf id} is a unique identifier for the element. If the argument is
            omitted then the current element is used, i.e. the one most recently
            created using inbuilt procedure {\bf element()} \autorefph{sec:ipelement}.
            
            {\bf prop} is the property identifier.
            
            {\bf l} is the value of the property. If no argument is passed then
            this procedure does nothing and returns. Note that if this procedure
            is called within another module or procedure and the property
            value is a parameter of that enclosing procedure then the inbuilt
            fuction {\bf nullarg()} \autorefph{sec:ifnullarg} may be used to enclose the
            argument to this procedure so that an omitted argument is ignored rather
            than the default parameter value of the enclosing procedure being used.
            
            See also - \\
            %{\bf elementLogProperty()} \autorefph{sec:lpelementLogProperty} for logical properties,
            {\bf elementIntProperty()} \autorefph{sec:lpelementIntProperty} for integer properties,
            {\bf elementBinProperty()} \autorefph{sec:lpelementBinProperty} for binary properties,
            {\bf elementHexProperty()} \autorefph{sec:lpelementHexProperty} for hexadecimal properties,
            {\bf elementFloatProperty()} \autorefph{sec:lpelementFloatProperty} for floating point properties,
            {\bf elementStrProperty()} \autorefph{sec:lpelementStrProperty} for string properties.






        \subsubsection{elementStrProperty - add a property of type "str" to an element}\label{sec:lpelementStrProperty}
            
            {\bf Library: xilinx/common.3pl}

            {\bf Prototype:} \\
            \vsp
            \begin{lstlisting}[gobble=12, language=threepl]
               global procedure
               elementStrProperty (str(id), str({prop, v}), varargs) {}
            \end{lstlisting}

            {\bf Output parameters:} \\
            none

            {\bf Input parameters:} \\            
            {\bf id} is the element identifier or null.\\
            {\bf prop} is the property identifier.\\
            {\bf v} is the property value.\\
            {\bf varargs} matches a number of string arguments listing the
                allowed property values.\\
        
            {\bf In-line target execution:} \\
            no

            {\bf description:} \\
            This procedure adds a property of type "str" to an element.
            {\bf id} is a unique identifier for the element. If the argument is
            omitted then the current element is used, i.e. the one most recently
            created using inbuilt procedure {\bf element()} \autorefph{sec:ipelement}.
            
            {\bf prop} is the property identifier.
            
            {\bf v} is the value of the property. If no argument is passed then
            this procedure does nothing and returns. Note that if this procedure
            is called within another module or procedure and the property
            value is a parameter of that enclosing procedure then the inbuilt
            fuction {\bf nullarg()} \autorefph{sec:ifnullarg} may be used to enclose the
            argument to this procedure so that an omitted argument is ignored rather
            than the default parameter value of the enclosing procedure being used.
            
            Optional arguments following {\bf prop} are strings specifying the allowed
            values for the property. If the property value is not identical to one
            of these strings then a fatal error will be forced. If the optional
            arguments are not supplied then the property will be defined
            unconditionally.
            
            See also - \\
            {\bf elementLogProperty()} \autorefph{sec:lpelementLogProperty} for logical properties,
            {\bf elementIntProperty()} \autorefph{sec:lpelementIntProperty} for integer properties,
            {\bf elementBinProperty()} \autorefph{sec:lpelementBinProperty} for binary properties,
            {\bf elementHexProperty()} \autorefph{sec:lpelementHexProperty} for hexadecimal properties,
            {\bf elementFloatProperty()} \autorefph{sec:lpelementFloatProperty} for floating point properties,
            %{\bf elementStrProperty()} \autorefph{sec:lpelementStrProperty} for string properties.

  



        \subsubsection{fdre - Create an FDRE flip-flop}\label{sec:lpfdre}
            
            {\bf Library: xilinx/common.3pl}

            {\bf Prototype:} \\
            \vsp
            \begin{lstlisting}[gobble=12, language=threepl]
               global procedure
               fdre (
                   value(q, "log")
                   <-
                   value(d, "log", false),
                   clock(c),
                   value(ce, "log", true),
                   value(r, "log", false),
                   str(init, "R")
                ) {
            \end{lstlisting}

            {\bf Output parameters:} \\
            {\bf q} is the flip-flop output.

            {\bf Input parameters:} \\            
            {\bf d} is the data input, default low.\\
            {\bf c} is the clock input.\\
            {\bf ce} is the optional clock enable input, default high.\\
            {\bf r} is the optional synchronous reset input, default low. \\
            {\bf init} is the initial value, default reset.
        
            {\bf In-line target execution:} \\
            no

            {\bf description:} \\
            This procedure creates a D-type flip-flop with clock enable and
            synchronous reset. The {\bf d}, {\bf ce} and {\bf r} inputs must
            have the same clock domain as the input clock. The {\bf q} output
            is set to the same clock domain.
            
            See also - \\
            {\bf fdse()} \autorefph{sec:lpfdse}\\
            {\bf fdrse()} \autorefph{sec:lpfdrse}\\
            {\bf fdce()} \autorefph{sec:lpfdce}\\
            {\bf fdpe()} \autorefph{sec:lpfdpe}.


        \subsubsection{fdrse - Create an FDRSE flip-flop}\label{sec:lpfdrse}
            
            {\bf Library: xilinx/common.3pl}

            {\bf Prototype:} \\
            \vsp
            \begin{lstlisting}[gobble=12, language=threepl]
               global procedure
               fdrse (
                   value(q, "log")
                   <-
                   value(d, "log", false),
                   clock(c),
                   value(ce, "log", true),
                   value(r, "log", false),
                   value(s, "log", false),
                   str(init, "R")
               ) {
            \end{lstlisting}

            {\bf Output parameters:} \\
            {\bf q} is the flip-flop output.

            {\bf Input parameters:} \\            
            {\bf d} is the data input, default low.\\
            {\bf c} is the clock input.\\
            {\bf ce} is the optional clock enable input, default high.\\
            {\bf r} is the optional synchronous reset input, default low. \\
            {\bf s} is the optional synchronous set input, default low. \\
            {\bf init} is the initial value, default reset.
        
            {\bf In-line target execution:} \\
            no

            {\bf description:} \\
            This procedure creates a D-type flip-flop with clock enable and
            synchronous reset and set. The {\bf d}, {\bf ce}, {\bf r} and {\bf
            s} inputs must have the same clock domain as the input clock. The
            {\bf q} output is set to the same clock domain.
            
            See also - \\
            {\bf fdre()} \autorefph{sec:lpfdre}\\
            {\bf fdse()} \autorefph{sec:lpfdse}\\
            {\bf fdce()} \autorefph{sec:lpfdce}\\
            {\bf fdpe()} \autorefph{sec:lpfdpe}.




        \subsubsection{fdse - Create an FDSE flip-flop}\label{sec:lpfdse}
            
            {\bf Library: xilinx/common.3pl}

            {\bf Prototype:} \\
            \vsp
            \begin{lstlisting}[gobble=12, language=threepl]
               global procedure
               fdse (
                   value(q, "log")
                   <-
                   value(d, "log", false),
                   clock(c),
                   value(ce, "log", true),
                   value(s, "log", false),
                   str(init, "S")
               ) {
            \end{lstlisting}

            {\bf Output parameters:} \\
            {\bf q} is the flip-flop output.

            {\bf Input parameters:} \\            
            {\bf d} is the data input, default low.\\
            {\bf c} is the clock input.\\
            {\bf ce} is the optional clock enable input, default high.\\
            {\bf s} is the optional synchronous set input, default low. \\
            {\bf init} is the initial value, default set.
        
            {\bf In-line target execution:} \\
            no

            {\bf description:} \\
            This procedure creates a D-type flip-flop with clock enable and
            synchronous set. The {\bf d}, {\bf ce} and {\bf s} inputs must have
            the same clock domain as the input clock. The {\bf q} output is set
            to the same clock domain.
            
            See also - \\
            {\bf fdre()} \autorefph{sec:lpfdre}\\
            {\bf fdrse()} \autorefph{sec:lpfdrse}\\
            {\bf fdce()} \autorefph{sec:lpfdce}\\
            {\bf fdpe()} \autorefph{sec:lpfdpe}.


        \subsubsection{fdce - Create an FDCE flip-flop}\label{sec:lpfdce}
            
            {\bf Library: xilinx/common.3pl}

            {\bf Prototype:} \\
            \vsp
            \begin{lstlisting}[gobble=12, language=threepl]
               global procedure
               fdce (
                   value(q, "log")
                   <-
                   value(d, "log", false),
                   clock(c),
                   value(ce, "log", true),
                   value(clr, "log", false),
                   str(init, "R")
               ) {
            \end{lstlisting}

            {\bf Output parameters:} \\
            {\bf q} is the flip-flop output.

            {\bf Input parameters:} \\            
            {\bf d} is the data input, default low.\\
            {\bf c} is the clock input.\\
            {\bf ce} is the optional clock enable input, default high.\\
            {\bf clr} is the optional asynchronous clear input, default low. \\
            {\bf init} is the initial value, default reset.
        
            {\bf In-line target execution:} \\
            no

            {\bf description:} \\
            This procedure creates a D-type flip-flop with clock enable and
            asynchronous clear. The {\bf d}, {\bf ce} and {\bf clr} inputs must
            have the same clock domain as the input clock. The {\bf q} output
            is set to the same clock domain.
            
            See also - \\
            {\bf fdre()} \autorefph{sec:lpfdre}\\
            {\bf fdse()} \autorefph{sec:lpfdse}\\
            {\bf fdrse()} \autorefph{sec:lpfdrse}\\
            {\bf fdpe()} \autorefph{sec:lpfdpe}.


        \subsubsection{fdpe - Create an FDPE flip-flop}\label{sec:lpfdpe}
            
            {\bf Library: xilinx/common.3pl}

            {\bf Prototype:} \\
            \vsp
            \begin{lstlisting}[gobble=12, language=threepl]
               global procedure
               fdpe (
                   value(q, "log")
                   <-
                   value(d, "log", false),
                   clock(c),
                   value(ce, "log", true),
                   value(pre, "log", false),
                   str(init, "S")
               ) {
            \end{lstlisting}

            {\bf Output parameters:} \\
            {\bf q} is the flip-flop output.

            {\bf Input parameters:} \\            
            {\bf d} is the data input, default low.\\
            {\bf c} is the clock input.\\
            {\bf ce} is the optional clock enable input, default high.\\
            {\bf pre} is the optional asynchronous set input, default low. \\
            {\bf init} is the initial value, default set.
        
            {\bf In-line target execution:} \\
            no

            {\bf description:} \\
            This procedure creates a D-type flip-flop with clock enable and
            asynchronous set. The {\bf d}, {\bf ce} and {\bf pre} inputs must
            have the same clock domain as the input clock. The {\bf q} output
            is set to the same clock domain.
            
            See also - \\
            {\bf fdre()} \autorefph{sec:lpfdre}\\
            {\bf fdse()} \autorefph{sec:lpfdse}\\
            {\bf fdrse()} \autorefph{sec:lpfdrse}\\
            {\bf fdce()} \autorefph{sec:lpfdce}.




        \subsubsection{generateConstraints - generate constraint file entries from I/O mode variable attributes}\label{sec:lpgenerateconstraints}
            
            {\bf Library: xilinx/common.3pl}

            {\bf Prototype:} \\
            \vsp
            \begin{lstlisting}[gobble=12, language=threepl]
               global procedure
               generateConstraints () {}
            \end{lstlisting}

            {\bf Output parameters:} \\
            none

            {\bf Input parameters:} \\            
            none
        
            {\bf In-line target execution:} \\
            no

            {\bf description:} \\
            
            This procedure traverses all input, output and clock mode variables
            and from their attribute maps generates constraint file entries. The
            procedure obtains a list of pointers to these variables by calling
            inbuilt function {\bf varlist()} \autorefph{sec:ifvarlist}. This
            procedure must be called once when there are no further calls to
            input(), output() or clock(). It is best to include the call to it in
            a late-order trailer module.




            

        \subsubsection{iopins - create FPGA pin specifications}\label{sec:lpiopins}
            
            {\bf Library: xilinx/common.3pl}

            {\bf Prototype:} \\
            \vsp
            \begin{lstlisting}[gobble=12, language=threepl]
               global procedure
               iopins(ident(id), type(t), var(pins, t), varargs) {}
            \end{lstlisting}

            {\bf Output parameters:} \\
            none

            {\bf Input parameters:} \\            
            {\bf id} is the identifier for a string or array of strings to be
            created.\\
            {\bf t} is the type of the following 'pins' argument.\\
            {\bf pins} is a collection of strings specifying FPGA pins.
        
            {\bf In-line target execution:} \\
            no

            {\bf description:} \\
            This procedure attaches attributes to a collection of FPGA pin
            strings. Optionally the procedure can create a variable containing
            the pin strings.
            
            For each pin string a map of attributes is created. This map is itself
            contained in a single map whose key is the pin string and value the
            attribute map for that pin.
            
            Parameter {\bf id} is an optional identifier which if given an argument
            will be used to create an immediate variable whose value is the collection
            of strings designating the FPGA pins.
            
            Parameter {\bf t} is the type of the pin argument. The argument may be
            type {\bf type} or it may be a type string. The type may be simply "str"
            for one pin, or more usually nested arrays and structs whose primitive
            types are only "str".
            
            Parameter {\bf pins} is the collection of pin strings, the argument to
            which may be either a compound constant or a variable of the designated
            type.
            
            Following these parameters a number of additional arguments may be
            provided to specify attributes associated with the pins. These additional
            arguments may be keymatch, i.e. attribute=value, or may be simple arguments.
            Where an argument is not keymatch it must be type "map" whose keys are
            attributes (strings) and values are attribute values.
            
            The arguments are processed left to right and where an attribute appears
            more than once the right-most value will prevail.
            
            Each attribute value must be a single value which will be added
            to each pin map.
            
            This procedure may be called multiple times, perhaps with different
            but overlapping collections of pins. Each time attributes will be
            attached to the attribute map of each pin. Where attribute entries
            already exist they will be overwritten. The variable identifier must
            be different in each case or may be omitted.
            
            It is important to note that the prime function of iopins() is to
            establish attributes associated with pins. As an optional byproduct a
            variable can be created which contains the collection of pins and this
            would be used to pass locations to library procedures input(), output()
            or clock() for external connections. The pin locations can however be
            passed to these procedures quite independently of iopins since only the
            pin strings are needed as the attributes are extracted from the
            attribute maps associated with the pins. It is of course usual to use
            the variable optionally created by iopins() to pass the location pins to
            the I/O procedures.
            
            The pin attributes established by iopins() are used by library
            procedures input(), output() and clock() but attributes can also be
            passed directly to these procedures as arguments to them in a similar
            way. It is a matter of choice whether to associate attributes directly
            with pins or instead to pass them directly to the input and output
            procedures. This may be partly decided by whether pins are connected
            to fixed external components or else to hardware which varies with the
            application.
            
            Attributes passed directly to input(), output() or clock() will
            override any attributes passed via iopins(). An attribute associated
            with a pin via iopins() that is passed to input(), output() or
            clock() must not conflict with that attribute on another pin in that
            call, though that attribute need not be on all pins.
            
            The attributes attached to pins by iopins() may be unused. Attributes
            are only used where pins are processed by library procedures input(),
            output() or clock() and then only those relevant to that procedure.
            Unrecognised attributes will result in a fatal error.
            
            The following table lists the allowed attributes that are recognised
            by iopins().
                    These attributes are not case-sensitive but will be converted to
        lower case.
         
        \btabularx
        \hline
        \multicolumn{4}{|l|}{\bf iopins()} \\
        \hline
        differential       & netlist    & log  & differential                      \\
        diff\_term         & constraint & log  & differential input termination    \\
        drive              & constraint & int  & drive current                     \\
        dutycycle          & 3PL        & float& clock duty cycle                  \\
        frequency          & 3PL        & float& clock frequency                   \\
        ibuf\_delay\_value & constraint & log  & input buffer delay                \\
        ibufg              & netlist    & log  & use an IBUFG instead of an IBUF   \\
        ifd\_delay\_value  & constraint & int  & input flip-flop delay             \\
        iostandard         & constraint & str  & voltage standard                  \\
        noshadow           & 3PL        & log  & don't duplicate output flip-flop  \\
        period             & 3PL        & float& clock period                      \\
        pulldown           & constraint & log  & pulldown                          \\
        pullup             & constraint & log  & pullup                            \\
        slew               & constraint & str  & slew rate                         \\
        tnm                & constraint & str  & timing name                       \\
        \hline
        \btabularxend

        {\bf diff\_term} if true indicates that a $100\Omega$ termination
        resistance be inserted between differential inputs
        
        {\bf differential} indicates that the pads are differential.
        Two locations must be supplied for each data bit.
        
        The {\bf drive} attribute is an integer which specifies the output drive
        current in mA for Xilinx output buffers (see Xilinx documentation).
        
        {\bf dutycycle} specifies the duty cycle of a clock input.
        
        {\bf frequency} specifies the frequency of a clock input.
        
        {\bf ibufdelay} gives the delay, 0 to 16, to be inserted in an input
        buffer path. The value has no particular units.
        
        {\bf ibufg} specifies that an IBUFG be used instead of an IBUF.
        
        {\bf ifddelay} gives the delay, 0 to 8, to be inserted in an input
        flip-flop input path. The value has no particular units.
        
        {\bf iostandard} is a string specifying one of the allowed
        I/O standards (see Xilinx documentation).
        
        {\bf noshadow} when true inhibits the duplication of flip-flops
        whose output is connected to an output buffer. By default any
        flip-flop whose output is connected to an output buffer is placed
        in the IO pad containing the buffer. If the output of the flip-flop
        is used elsewhere then the flip-flop is duplicated in a CLB with
        all inputs in parallel with the pad flip-flop and the output of the
        duplicate used for internal logic. If {\bf noshadow} is true then
        duplication will not be done, the flip-flop will of necessity
        reside in a CLB and its output will be routed to the output buffer.
        This will result in much worse clock-to-output time. Setting this
        attribute ensures that the pad output genuinely reflects the
        flip-flop state and not that of a duplicate but at the cost of
        delayed output.
        
        {\bf period} specifies the period of a clock input.
        
        If {\bf pullup} or {\bf pulldown} are true the output will be given
        a weak pullup or pulldown. It is not allowed for both these
        attributes to be true.
        
        {\bf pullup}
        
        The {\bf slew} attribute is the string "fast" or "slow" and selects
        one of two slew rates for Xilinx output buffers (see Xilinx documentation).
        
        {\bf tnm} is the name of a timing group in which the output pads
        are to be included for timing constraint purposes.
        

                        
            Typical useage might be -
            \begin{lstlisting}[gobble=12, language=threepl]
               static(led, "bits:4");
               var(tnms, "[2]str", {"l1", "l2"});
               iopins(led_pins, "[4]str", {"A1", "B1", "A2", "C2"}, iostandard="lvcmos33", pullup=true);
               iopins(, "[2]str", {"B1", "A2"}, tnm=tnms);
               output(,, led, loc=led_pins, drive=24);
            \end{lstlisting}
            which will create a variable {\bf led\_pins} in the current scope
            which is an array of 4 strings. The I/O standard "lvcmos33" will be
            established for each of the four pins (inputs or outputs) and a
            pullup resistor is specified for each. For two of the pins separate
            timing name attributes have been added.
            
            The following example illustrates that attribute values for pins
            cab be overridden by an additional call to {\bf iopins()} -
            \begin{lstlisting}[gobble=12, language=threepl]
               static(led, "bits:4");
               iopins(led_pins, "[4]str", {"A1", "B1", "A2", "C2"}, pullup=true);
               iopins(, "[2]str", {"B1", "A2"}, pullup=false);
               output(,, led, loc=led_pins);
            \end{lstlisting}
            
            An alternative is to pass attribute values directly through calls to
            {\bf input()}, {\bf output}, or {\bf clock()} which will override attrbute values
            previously specified via {\bf iopins()}, e.g.
            \begin{lstlisting}[gobble=12, language=threepl]
               static(led, "bits:4");
               iopins(led_pins, "[4]str", {"A1", "B1", "A2", "C2"}, pullup=true);
               iopins(, "[2]str", {"B1", "A2"}, pullup=false);
               output(,, led, loc=led_pins, pullup=true);
            \end{lstlisting}
            
            
            This procedure is used in conjunction with inbuilt procedures
            {\bf input()} \autorefph{sec:ipinput}, \\
            {\bf output()} \autorefph{sec:ipoutput} and \\
            {\bf clock()} \autorefph{sec:ipclock}.









        \subsubsection{maxdelay - specify max delay between timing groups}\label{sec:lpmaxdelay}
            \index{constraint!timing}
            \index{timing constraint}
            
            {\bf Library: xilinx/common.3pl}

            {\bf Prototype:} \\
            \vsp
            \begin{lstlisting}[gobble=12, language=threepl]
               global procedure
               maxdelay (str({group1, group2}), float(time)) {}
            \end{lstlisting}

            {\bf Output parameters:} \\
            none

            {\bf Input parameters:} \\            
            {\bf group1} is a string giving the timing name of a netlist group.\\
            {\bf group2} is a string giving the timing name of a netlist group.\\
            {\bf time} is the maximum delay in second allowed
            for signals propagating from signal sources in group1 to
            signal destinations in group2.
        
            {\bf In-line target execution:} \\
            no

            {\bf description:} \\
            This procedure sets up the timing constraints for paths between
            components in different timing groups.
            
            It will generate a {\bf TIMESPEC} constraint in the ISE NCF file or
            a corresponding {\bf set\_max\_delay} constraint in the Vivado XDC
            file.

            For input delay constraints -\\
            
            {\bf group1} is timing group name for
            the input ports and {\bf group2} is the timing name for the clock.
            
            For ISE {\bf time} specifies the maximum delay in ns for all signal
            paths from all the relevant input ports to elements clocked by the
            clock. This is a maximum delay time constraint from port to clock
            internal to the FPGA. This maximum delay must be met by ISE
            to satisfy the constraint.
            
            For Vivado {\bf time} specifies the maximum delay in the external
            signals from the clock edge to arrival at the input ports.
            This is the actual maximum delay time from external clock edge to
            input ports external to the FPGA where the notional external clock
            is phase aligned with the internal clock. If the actual external
            clock is skewed relative to the internal clock then this skew
            must be adjusted for in the specified time. This delay time is
            used by Vivado along with the clock period to determine the
            maximum time delay between pads and clocked elements that must
            be achieved during place-and-route.


            For output delay constraints {\bf group1} is timing name for the
            clock and {\bf group2} is the timing group name for the output
            ports.
            
            For ISE {\bf time} specifies the maximum delay in ns for all signal
            paths from elements clocked by the clock to all the relevant output
            ports. This is a maximum delay time constraint from clock to port
            internal to the FPGA. This maximum delay must be met by ISE to
            satisfy the constraint.
            
            For Vivado {\bf time} specifies the maximum delay in the external
            signals from the output ports to the external clock edge . This is
            the actual maximum delay time from output ports to external clock
            edge external to the FPGA where the notional external clock is phase
            aligned with the internal clock. If the actual external clock is
            skewed relative to the internal clock then this skew must be
            adjusted for in the specified time. This delay time is used by
            Vivado along with the clock period to determine the maximum time
            delay between clocked elements and pads that must be achieved during
            place-and-route.
            
            {\bf maxdelay()} only adds a struct to a list of delay structs
            in common.3pl. A later call to generateConstraints() builds
            the constraint elements for ISE or Vivado as appropriate.

            See also - \\         
            {\bf addncfline()} \autorefph{sec:lpaddncfline}\\
            {\bf append\_ncf\_file()} \autorefph{sec:lpappncf}\\
            {\bf elementaddncf()} \autorefph{sec:lpelementaddncf}\\
            {\bf addxdcline()} \autorefph{sec:lpaddxdcline}\\
            {\bf period\_mul()} \autorefph{sec:lpperiodmul}\\
            {\bf period\_div()} \autorefph{sec:lpperiodmul}\\
            {\bf netaddncf()} \autorefph{sec:lpnetaddncf}.





        \subsubsection{netaddncf - add constraint to NCF file for a net}\label{sec:lpnetaddncf}
            \index{constraint}
            
            {\bf Library: xilinx/common.3pl}

            {\bf Prototype:} \\
            \vsp
            \begin{lstlisting}[gobble=12, language=threepl]
               global procedure
               addncf (str(s), value(sig)) {}
               netaddncf (value(sig), str({attr, val})) {}
            \end{lstlisting}

            {\bf Output parameters:} \\
            none

            {\bf Input parameters:} \\            
            {\bf sig} is .\\
            {\bf attr} is .\\
            {\bf val} is .
        
            {\bf In-line target execution:} \\
            no

            {\bf description:} \\
            An attribute or constraint is applied to a net. Argument
            {/bf sig} is a target variable that identifies the subject net. The
            attribute/constraint {/bf attr} and the value {/bf val} are both
            strings.
            
            See also - \\
            {\bf addncfline()} \autorefph{sec:lpaddncfline}\\
            {\bf append\_ncf\_file()} \autorefph{sec:lpappncf}\\
            {\bf elementaddncf()} \autorefph{sec:lpelementaddncf}\\
            {\bf addxdcline()} \autorefph{sec:lpaddxdcline}\\
            {\bf maxdelay()} \autorefph{sec:lpmaxdelay}\\
            {\bf period\_mul()} \autorefph{sec:lpperiodmul}\\
            {\bf period\_div()} \autorefph{sec:lpperiodmul}.









        \subsubsection{period\_div - timing group for divided clock}\label{sec:lpperioddiv}
            
            {\bf Library: xilinx/common.3pl}

            {\bf Prototype:} \\
            \vsp
            \begin{lstlisting}[gobble=12, language=threepl]
               global procedure
               period_div (str(newgroup, group), float(ratio)) {}
            \end{lstlisting}

            {\bf Output parameters:} \\
            none

            {\bf Input parameters:} \\            
            {\bf newgroup} is a new timing group identifier. \\
            {\bf group} is a timing group identifier. \\
            {\bf ratio} is a period ratio. \\
        
            {\bf In-line target execution:} \\
            no

            {\bf description:} \\
            This procedure creates a new timing group whose name is
            given by {\bf newgroup}. The period of the new group is
            the period of the old timing group {\bf group} divided
            by {\bf ratio}.
            
            See also - \\
            {\bf addncfline()} \autorefph{sec:lpaddncfline}\\
            {\bf append\_ncf\_file()} \autorefph{sec:lpappncf}\\
            {\bf elementaddncf()} \autorefph{sec:lpelementaddncf}\\
            {\bf addxdcline()} \autorefph{sec:lpaddxdcline}\\
            {\bf maxdelay()} \autorefph{sec:lpmaxdelay}\\
            {\bf period\_mul()} \autorefph{sec:lpperiodmul}\\
            {\bf netaddncf()} \autorefph{sec:lpnetaddncf}.




        \subsubsection{period\_mul - timing group for divided clock}\label{sec:lpperiodmul}
            
            {\bf Library: xilinx/common.3pl}

            {\bf Prototype:} \\
            \vsp
            \begin{lstlisting}[gobble=12, language=threepl]
               global procedure
               period\_mul (str(newgroup, group), float(ratio)) {}
            \end{lstlisting}

            {\bf Output parameters:} \\
            none

            {\bf Input parameters:} \\            
            {\bf newgroup} is a new timing group identifier. \\
            {\bf group} is a timing group identifier. \\
            {\bf ratio} is a period ratio. \\
        
            {\bf In-line target execution:} \\
            no

            {\bf description:} \\
            This procedure creates a new timing group whose name is
            given by {\bf newgroup}. The period of the new group is
            the period of the old timing group {\bf group} multiplied
            by {\bf ratio}.
            
            See also - \\
            {\bf addncfline()} \autorefph{sec:lpaddncfline}\\
            {\bf append\_ncf\_file()} \autorefph{sec:lpappncf}\\
            {\bf elementaddncf()} \autorefph{sec:lpelementaddncf}\\
            {\bf addxdcline()} \autorefph{sec:lpaddxdcline}\\
            {\bf maxdelay()} \autorefph{sec:lpmaxdelay}\\
            {\bf period\_div()} \autorefph{sec:lpperiodmul}\\
            {\bf netaddncf()} \autorefph{sec:lpnetaddncf}.
 

        \subsubsection{processPinAttributes - process pin attributes}\label{sec:lpprocPinAtt}
            
            {\bf Library: xilinx/common.3pl}

            {\bf Prototype:} \\
            \vsp
            \begin{lstlisting}[gobble=12, language=threepl]
               global procedure
               processPinAttributes (var(ptr, "->")) {}
            \end{lstlisting}

            {\bf Output parameters:} \\
            none

            {\bf Input parameters:} \\            
            {\bf ptr} is a pointer to an input, output or clock mode variable.
        
            {\bf In-line target execution:} \\
            no

            {\bf description:} \\
            
            This procedure is called following the creation of every input, output
            or clock mode variable. This procedure must have been previously
            registered with the 3PL system by calling
            
            \begin{lstlisting}[gobble=12, language=threepl]
               attributeprocedure(processPinAttributes);
            \end{lstlisting}
            
            prior to any calls to inbuilt procedures input(), output() or clock().
            This is best done by placing the call in a leader module, ordered
            within leader modules such that it occurs early enough to precede the
            above variable creation calls. 
        
            The argument is a pointer to an input, output or clock mode variable.
            This procedure accesses the attribute map of the variable concerned and
            extracts the 'loc' attribute if present. From the 'loc' attribute it
            consults global map {\bf pinAttributes} which has been created by
            library procedure {\bf iopins()} \autorefph{sec:lpiopins} using each pin
            (string) as a key. The value from that map is another map containing
            attributes for that pin. These attributes are added to the attributes
            map of the variable if they are not already present since attributes
            passed as arguments to the variable creation procedure should override
            attributes derived from pins.
            
            When an attribute is extracted from pin attribute maps it must be added
            to the input, output or clock variable attribute map as an array with
            one entry per pin since an attribute may have different values amongst
            the pins.           
            
            The procedure also expands most attributes from the variable attribute
            map to arrays if they are single values (some attributes, such as
            frequency or readdomain, are required to remain single valued as they apply
            to the variable itself, not the input/output buffers or pads).
            
            In the case of an input clock with differential input attributes
            attached to the pair of pins are redundant since the clock attributes
            except {\bf loc} are required to be single values. It is reasonable
            to expect that attributes associated with pins will be a pair of identical
            values and processPinAttributes() may use either (it will actually
            extract them in turn and the 2nd one will prevail).
            
            Note that processPinAttributes() ignores attributes that it does not
            recognise. At present iopins() has a list of attributes that it will
            accept, which it checks for type. It will give a fatal error on an
            unrecognised attribute (though input(), output() and clock() will
            accept unrecognised attributes for addition to their attribute maps
            without comment). This may be relaxed in the future so that additional
            arbitrary attributes can be added for software use.
            
            For information on attributes see {\bf input()} \autorefph{sec:ipinput},
            {\bf output()} \autorefph{sec:ipoutput} or {\bf clock()}
            \autorefph{sec:ipclock}.
            
            See also {\bf attributeprocedure()} \autorefph{sec:ipattributeproc}.



        \subsubsection{pullup - connect pullups to a bus}\label{sec:lppullup}
            
            {\bf Library: xilinx/common.3pl}

            {\bf Prototype:} \\
            \vsp
            \begin{lstlisting}[gobble=12, language=threepl]
               global procedure
               pullup (value(in)) {}
            \end{lstlisting}

            {\bf Output parameters:} \\
            none

            {\bf Input parameters:} \\            
            {\bf in} is a selectvalue mode variable. \\
        
            {\bf In-line target execution:} \\
            no

            {\bf description:} \\
            This procedure connects pullup resistors to a selectvalue
            variable (an internal bus).


            

        \subsubsection{setiostandard - set the I/O voltage standard for a pin or pins}\label{sec:lpsetiostandard}
            
            {\bf Library: xilinx/common.3pl}

            {\bf Prototype:} \\
            \vsp
            \begin{lstlisting}[gobble=12, language=threepl]
               global procedure
               setiostandard (str(iostandard), varargs) {}
            \end{lstlisting}

            {\bf Output parameters:} \\
            none

            {\bf Input parameters:} \\            
            {\bf iostandard} is a string giving the I/O voltage standard.\\
            For second argument see description.
        
            {\bf In-line target execution:} \\
            no

            {\bf description:} \\
            This procedure established I/O voltage standards for one or more FPGA
            inputs or outputs. The first argument is a string giving the voltage
            standard. The second argument is a string or a string array of any
            number of dimensions, the strings being the FPGA pin designations. The
            voltage standard will be established for the specified pins (it is
            stored in a map where each pin is a key and the associated value is the
            standard string). This procedure may be called again to change the
            standard established for pins on a previous call since the I/O
            standards are only used at the completion of immediate execution.
            
            This procedure is used in conjunction with library function {\bf
            getiostandard()} \autorefph{sec:lfgetiostandard} and library procedure
            {\bf iopins()} \autorefph{sec:lpiopins}.












    \subsection{library functions}




        \subsubsection{clockused - determine if a clock variable has been used}\label{sec:lfclockused}
            
            {\bf Library: xilinx/common.3pl}

            {\bf Prototype:} \\
            \vsp
            \begin{lstlisting}[gobble=12, language=threepl]
               global function
               clockused (c) {}
            \end{lstlisting}

            {\bf Return mode and type:} \\
            immediate, log

            {\bf Input parameters:} \\            
            {\bf c} is a clock variable

            {\bf description:} \\
            This function returns true if the clock variable passed as argument
            has been used to drive logic.
            
            The function can be used to avoid generating timing constraints, and
            even the clock generator itself, if a clock variable is not used
            to drive logic. This determination can only be made at the completion
            of application code and so would normally be done in a trailer
            module.




        \subsubsection{del - variable length delay line}\label{sec:lfdel}
            
            {\bf Library: xilinx/common.3pl}

            {\bf Prototype:} \\
            \vsp
            \begin{lstlisting}[gobble=12, language=threepl]
               global function
               del (
                   value(in),
                   int(len, 1),
                   var(init, "[]int"),
                   value(step, "log", true),
                   value(varlen, "uint:4"),
                   value(reset, "log"),
               ) {}
            \end{lstlisting}

            {\bf Return mode and type:} \\
            value, type(in)

            {\bf Input parameters:} \\            
            {\bf in} is the input value. It may be any target type.\\
            {\bf len} is the length of the delay line.\\
            {\bf init} is an optional initialisation.\\
            {\bf step} is an optional step control.\\
            {\bf varlen} is an optional variable length\\
            {\bf reset} is an optional reset.

            {\bf description:} \\
            This function is a delay line implemented using CLB shift registers.
            On each clock cycle the input is written to the line and the
            contents move one item right. The rightmost item appears as the out
            value. This calls inbuilt module {\bf del()} \autorefph{sec:imdel}.
            This function can operate as a sequence generator if the output is
            connected back to the input.

            This has up to 6 input arguments. Input arguments {\bf len} and {\bf
            init} are immediate mode. {\bf step} may be immediate or target mode
            as described below. The other input arguments must not be modes {\bf
            immediate}, {\bf queue}, {\bf memory} or {\bf clock} and value mode
            arguments must not contain queue reads.
            
            The argument to {\bf step} must be type "log" and may be as follows -
            \blist
            \item   no argument - the delay will advance on each clock cycle
            \item   an immediate mode value (or constant value mode) which is
                    {\bf true} - the delay will advance on each clock cycle
            \item   an immediate mode value (or constant value mode) which is
                    {\bf false} - the delay will advance on clock cycles in
                    which the function is evaluated in the target code, e.g.
                    on used in assignment
            \item   a static, selectvalue or value mode value - the delay will
                    advance on clock cycles when {\bf step} is true.
            \blend
            
            For example -
            \begin{lstlisting}[gobble=12, language=threepl]
                when (l)
                    a <- del(b, 10, false);
            \end{lstlisting}
            will shift the current value of {\bf b} into the delay line every
            time the assignment to {\bf a} is made ({\bf l} is true), the
            assigned value being popped off the end.
            
            Whereas -
            \begin{lstlisting}[gobble=12, language=threepl]
                when (l)
                    a <- del(b, 10);
            \end{lstlisting}
            will shift {\bf b} into the delay line every clock cycle, popping
            off the vaue at the other end. When {\bf l} is true the value popped
            off will also be assigned to {\bf a}.
            \begin{lstlisting}[gobble=12, language=threepl]
                when (l)
                    a <- del(b, 10, true);
            \end{lstlisting}
            will behave the same way.

            Where {\bf step} is variable during target execution the shifting
            is controlled solely by step while the assignment, controlled by
            {\bf l}, assigns the end value of the delay line to {\bf a}
            regardless of whether a shift has occured or not.
            \begin{lstlisting}[gobble=12, language=threepl]
                when (l)
                    a <- del(b, 10, step);
            \end{lstlisting}

            The limitation stated above that inputs cannot contain queue reads
            can be overcome in some cases. For example to use a queue as data input
            we can use the examine operator as long as a queue wait controls the
            delay step. For example to form a running average of length 8 -

            \begin{lstlisting}[gobble=12, language=threepl]
                queue(q, "uint:8");
                static(rsum, "uint:11");
                queue(rav, "uint:8");

                seq {
                    rsum <- rsum + q - del(?q, 8,, false);
                    rav <<- rsum / 8;
                }
            \end{lstlisting}
            The data input to the {\bf del()} function is a queue examine, not
            a queue read. The examine operator provides the value of the
            top of a queue regardless of the validity (e.g. it might be empty).
            This is fine of the statement only executes when we have separately
            determined that the queue is avaiable for reading, which is the case here.
            The {\bf del()} function as noted above will only step when the
            function is evaluated if the step argument is immediate constant
            {\bf false}. In this example the function will not be evaluated
            until {\bf q} is available to read as it is included in the RHS expression
            along with the function call. AT that point the function will be evaluated
            and the delay line will step.

            For large delays this uses a lot of CLBs and the alternative library
            function {\bf delay()} \autorefph{sec:lfdelay} in library stdlib.3pl
            might be preferable. That library function will call {\bf
            del()} in any case if the delay value is small.

            This has 1 to 6 input arguments and 1 output argument. Input
            arguments {\bf len} and {\bf init} are immediate mode. The other
            input arguments must not be modes {\bf immediate}, {\bf queue}, {\bf
            memory} or {\bf clock} and expression arguments must not contain
            queue reads.

            The 1st input argument, {\bf in}, is the delay line data input. It
            can be any target type.

            The optional input argument, {\bf len}, specifies a fixed delay.
            This must be an immediate int. If this argument is 0 and
            the 5th argument {\bf varlen} is also omitted, the output will
            simply be connected to the input to give zero delay.

            The optional input argument, {\bf init}, is an initialisation
            array. This  must be immediate mode and an array of the same type as
            {\bf in}. It may be smaller than the delay line length in which case
            the output end of the delay line will be zero initialised. The low
            order entries in the initialisation array are loaded in the input
            end of the delay line. If initialisation values are too large their
            high-order bits will be truncated.

            The optional input argument, {\bf step}, is a line step enable.
            This  must be type log. If this argument is used the line will only
            advance when {\bf step} is true, otherwise it will advance on every
            clock cycle.

            The optional input argument, {\bf varlen}, is a variable line
            length. This must be type uint. This argument may be limited or
            unavailable depending on the type of FPGA. For Xilinx the variable
            length is limited to values between 0 and 15 giving delays of 1 to
            16 respectively. A wider argument will be truncated and a narrower
            one zero padded. Note that when the line length is changed data may
            be skipped or repeated. If the line length is changed while stepping
            is inhibited, the line length argument is an address with which the
            contents of the line can be non-destructively examined.

            The optional input argument, {\bf reset}, returns the
            line to its initial state when true.

            The delay line data output is the value of the function which is
            the same type as the data input {\bf in}.

            If the delay line is initialised, these initial values appear at the
            output during the initial {\em len} line steps, where {\em len} is the
            fixed delay (or the initial delay selected by {\bf varlen}
            if used).

            For Xilinx the following apply -
            \blist
            \item   For FPGAs prior to Virtex5 the variable length argument is limited to values
                    between 0 and 15 inclusive, giving line lengths of 1 to
                    16 respectively. From Virtex5 onwards the lengths are 1 to 32.
            \item   If the reset is used, the variable length argument cannot be
                    used.
            \item   If the reset is used or the line length is 1, the delay line
                    will be implemented using flip-flops, otherwise it will be
                    implemented using CLB LUT shift registers.
            \blend
            Implementation using CLB LUT shift registers minimises resources,
            requiring only one CLB LUT for every 16-bit delay line.

            {\bf Note: the fixed delay parameter {\em len} gives the number of
            cycles to be delayed but the variable delay parameter {\em varlen}
            is the number of cycles delay minus one!}      

            In Xilinx devices there is a substantial delay between the clock
            edge and the arrival of the shift register signal at the output of
            the logic block (2.8nS to 3.5nS for a -6 speed XC2VP and about 1.15
            to 1.43nS for  a -2 speed XC5V). The delay for a flip-flop is very
            much smaller. For this reason a delay line with fixed length is
            implemented using a flip-flop for the last delay stage. If a
            variable length delay line is required, this flip-flop is not
            inserted and hence the significant output delay must be considered.
            It is preferable in that case to assign the delay line output
            directly to a variable, i.e. with no intervening combinatorial
            logic, to avoid introducing any extra path delay.

            See also library function {\bf delay()} \autorefph{sec:lfdelay}
            which constructs a fixed-length delay line but for long delays
            provides an option of using block RAM rather than shift registers.



        \subsubsection{elementcheck - check that a Xilinx element is allowed in the current device family}\label{sec:lfelementcheck}
            
            {\bf Library: xilinx/common.3pl}

            {\bf Prototype:} \\
            \vsp
            \begin{lstlisting}[gobble=12, language=threepl]
               global function
               elementCheck(str(element)) {}
            \end{lstlisting}

            {\bf Return mode and type:} \\
            immediate, log

            {\bf Input parameters:} \\            
            {\bf element} is the name of a Xilinx library element.

            {\bf description:} \\
            This function checks that the Xilinx library element named by the
            argument is available for the current Xilinx device family. Allowed elements
            for each device family are kept in an abbreviated map in xilinx/common.3pl,
            i.e. only a subset of elements are included, namely those used explicitly
            in 3PL libraries. The map can be expanded when required.
            
            The result of the check is returned, true if the element is available in the
            current device family.
            
            See also library procedure {\bf elementcheck()} \autorefph{sec:lpelementcheck}
            which will issue a fatal error message on failure, rather than returning
            the check result.


        \subsubsection{getiostandard - get the I/O standard for a pin or pins}\label{sec:lfgetiostandard}
            
            {\bf Library: xilinx/common.3pl}

            {\bf Prototype:} \\
            \vsp
            \begin{lstlisting}[gobble=12, language=threepl]
               global function
               getiostandard (varargs) {}
            \end{lstlisting}

            {\bf Return mode and type:} \\
            immediate, str

            {\bf Input parameters:} \\            
            See description.

            {\bf description:} \\
            This function returns a string giving the I/O standard for one
            or more FPGA inputs or outputs. The argument must be a string or a string array of any
            number of dimensions, the strings being the FPGA pin designations.
            Multiple pins must have the same I/O standard or a fatal error occurs.
            
            The I/O standard for one or more pins is established or changed by calling
            library procedure {\bf setiostandard()} \autorefph{sec:lpsetiostandard}.
            
            This function is used in conjunction with library procedures {\bf
            setiostandard()} \autorefph{sec:lpsetiostandard} and {\bf iopins()}
            \autorefph{sec:lpiopins}.





        \subsubsection{reverse - reverse a value}\label{sec:lfreverse}
            
            {\bf Library: xilinx/common.3pl}

            {\bf Prototype:} \\
            \vsp
            \begin{lstlisting}[gobble=12, language=threepl]
               global function
               reverse (value(in)) {}
            \end{lstlisting}

            {\bf Return mode and type:} \\
            value, type(in)

            {\bf Input parameters:} \\            
            {\bf in} is the operand.

            {\bf description:} \\
            This function returns input operand with bits reversed. Since this
            does a simple bit-wise reversal it only makes sense if the input
            type can be sensibly reversed.
           




        \subsubsection{sequence - generate a sequence}\label{sec:lfsequence}
            
            {\bf Library: xilinx/common.3pl}

            {\bf Prototype:} \\
            \vsp
            \begin{lstlisting}[gobble=12, language=threepl]
               global function
               sequence (var(s), type(stype), value(step, "log", true)) {}
            \end{lstlisting}

            {\bf Return mode and type:} \\
            value, "bits:"

            {\bf Input parameters:} \\            
            {\bf s} is the sequence array.\\
            {\bf stype} is the data type.\\
            {\bf step} is an optional input to control data movement.

            {\bf description:} \\
            This function generates a cyclic sequence of values, these
            appearing on the output value variable. The sequence is
            determined by the input argument {\bf s} which must be an array
            of immediate values. The type of these values is given by the 2nd
            argument, {\bf stype}, and may be of compound type. The 3rd input
            argument, {\bf step}, is optional and is used to control when the
            sequence advances, otherwise this occurs on every clock edge.
            The sequence is produced starting with subscript 0 of the sequence
            data and moving up the array.




    \subsection{library classes}
    
        None.








